% Vietnamese translation for OI Public Library
% Source-PDF: 数据结构 DATA STRUCTURE/Link Cut Trees.pdf
\documentclass[11pt]{article}
\usepackage{oipl-vn}

\title{Link-Cut Trees}
\author{}
\date{}

\begin{document}
\maketitle

\origtitle{Link-Cut Trees}
\sourcepath{data-structure/link-cut-trees.pdf}

Link-Cut Tree là cấu trúc dữ liệu do Sleator và Tarjan đề xuất để giải các bài
toán \term{cây động}.  Với cấu trúc này, mỗi thao tác cơ bản trên cây động có
thể thực hiện trong thời gian khấu hao $O(\log n)$.

Tài liệu này lần lượt giới thiệu khái niệm của Link-Cut Tree, các thao tác được
hỗ trợ, cách viết Splay Tree bên trong Link-Cut Tree, rồi đi qua một số bài
tập thực hành.

\section{Định nghĩa Link-Cut Tree}

Một vài thuật ngữ tiếng Anh thường gặp:
\begin{itemize}
  \item \term{access}: truy cập;
  \item \term{preferred}: ưu tiên, được chọn;
  \item \term{path}: đường đi;
  \item \term{auxiliary}: phụ trợ.
\end{itemize}

Link-Cut Tree phân rã một cây thành nhiều chuỗi trên cây, rồi lưu riêng từng
chuỗi trong một Splay Tree.  Khi cần thao tác trên một đường đi trong cây gốc,
ta dùng Splay để tách và ghép các chuỗi sao cho đường đi đó nằm trong cùng một
Splay Tree; sau đó chỉ cần thao tác trên Splay Tree ấy.

Các khái niệm chính là:
\begin{itemize}
  \item \term{Preferred Child}: con ưu tiên.
  \item \term{Preferred Edge}: cạnh ưu tiên.
  \item \term{Preferred Path}: đường ưu tiên.
  \item \term{Auxiliary Tree}: cây phụ trợ, tức Splay Tree lưu một đường ưu
  tiên.
  \item \term{Path Parent}: cha của đường.
\end{itemize}

Ví dụ, xét một cây gốc tại $1$.  Nếu con ưu tiên của $1$ là $2$, con ưu tiên
của $2$ là $4$, còn $3,4,5,6$ không có con ưu tiên, thì đường $1-2-4$ là một
đường ưu tiên.  Cũng giống phân rã nặng nhẹ, Link-Cut Tree biến một cây thành
nhiều chuỗi tuyến tính; điểm khác là mỗi chuỗi được lưu bằng Splay Tree để tiện
tách và ghép động.

Trong Splay Tree của một đường ưu tiên, đỉnh gần gốc cây gốc hơn nằm về phía
trái, đỉnh sâu hơn nằm về phía phải; nói cách khác, khóa sắp xếp là độ sâu
\texttt{deep}.  Gốc của Splay Tree có một con trỏ cha đặc biệt: đó là
\term{path parent}, bằng cha trong cây gốc của đỉnh nông nhất trên đường này.

Nếu xem mỗi đường ưu tiên như một đỉnh và xem các cạnh không ưu tiên như cạnh
giữa các đỉnh ấy, ta thu được một cây ảo.  Cây ảo này thường nhỏ hơn cây ban
đầu về cả chiều rộng lẫn chiều cao.  Nếu cấu trúc cây ảo không đổi, rồi trên
mỗi đường ta dùng một cấu trúc tuyến tính để tăng tốc, ta nhận được ý tưởng của
phân rã nặng nhẹ.  Có thể hiểu phân rã nặng nhẹ là một Link-Cut Tree tĩnh.  Tuy
vậy, điểm cốt lõi của Link-Cut Tree không chỉ là nén cây, mà là khả năng đưa
đường cần truy vấn thành một đường chính trong thời gian $O(\log n)$ khấu hao.

\section{Các thao tác cơ bản}

\subsection{\texttt{access}}

\texttt{access} là nền tảng của mọi thao tác trong Link-Cut Tree.  Sau khi gọi
\texttt{access(v)}, đường từ gốc của cây chứa $v$ tới $v$ trở thành một đường
ưu tiên mới.  Nói cách khác, tất cả các đỉnh trên đường từ gốc tới $v$ trở
thành con ưu tiên tương ứng, các cạnh trên đường ấy trở thành cạnh ưu tiên, và
đường đó được lưu trong cùng một Splay Tree.

Cách làm là lặp trên các cây phụ trợ.  Khi đang xử lý đỉnh $x$, trước hết
\texttt{splay(x)} để đưa $x$ lên gốc Splay Tree hiện tại.  Vì đỉnh vừa được
truy cập sẽ không giữ con ưu tiên cũ ở phía sâu hơn, ta cắt con phải của $x$.
Sau đó nối phần đường đã xử lý trước đó vào con phải của $x$.  Tiếp theo đi
lên \texttt{path parent} của Splay Tree hiện tại và lặp lại.  Cuối cùng ta thu
được một đường từ gốc cây gốc tới $x$.

Mô phỏng bằng lời cho \texttt{access(2)}: đưa $2$ lên gốc của Splay Tree chứa
nó và đặt con phải của $2$ bằng $0$; sau đó xử lý path parent, ví dụ là $3$,
đưa $3$ lên gốc, cắt con phải của $3$ và nối đường chứa $2$ vào đó.  Nếu path
parent kế tiếp là $1$, tiếp tục đưa $1$ lên gốc, cắt con phải của $1$ và nối
đường chứa $3$.  Sau bước cuối, đường chính là đường từ gốc đến $2$.

Lý do luôn thao tác với con phải là trong cây phụ trợ, phía phải chứa các đỉnh
sâu hơn.  Khi một đỉnh được đưa lên gốc Splay Tree, con trái của nó là đoạn
đường về phía gốc cây gốc, còn con phải là đoạn đi xuống sâu hơn; ta cần thay
đoạn đi xuống ấy bằng đường vừa truy cập.

\begingroup
\small
\begin{verbatim}
void access(int root)
{
    int temp = 0;
    while (root)
    {
        splay(root);
        t[root].ch[1] = temp;
        updata(root);
        temp = root;
        root = t[root].fa;
    }
}
\end{verbatim}
\endgroup

\subsection{\texttt{lca}}

Nếu không trực tiếp hỏi LCA thì thao tác này ít khi cần dùng; với cây tĩnh lại
càng không cần dùng Link-Cut Tree.  Nhưng nếu cây động hỏi trực tiếp LCA, ta có
thể dựa vào \texttt{access}.

Giả sử cần tìm LCA của $9$ và $6$.  Trước hết gọi \texttt{access(9)}, khi đó
đường từ gốc tới $9$ được dựng thành đường chính.  Sau đó gọi
\texttt{access(6)}.  Hai đường từ gốc tới $9$ và từ gốc tới $6$ có một đoạn
chung; đoạn chung này không cần thay đổi trong lần \texttt{access} thứ hai.
Vì vậy, đỉnh cuối cùng bị thay đổi con phải chính là LCA.  Trong Splay Tree,
đỉnh cuối cùng được truy cập cũng vừa được xoay lên gốc, nên ta chỉ cần trả về
biến tạm cuối cùng của \texttt{access}.

\begingroup
\small
\begin{verbatim}
int access(int root)
{
    int temp = 0;
    while (root)
    {
        splay(root);
        t[root].ch[1] = temp;
        temp = root;
        root = t[root].fa;
    }
    return temp;
}

int lca(int root1, int root2)
{
    access(root1);
    return access(root2);
}
\end{verbatim}
\endgroup

Trong ví dụ của tài liệu gốc, sau khi thực hiện \texttt{access(9)} rồi
\texttt{access(6)}, đỉnh $5$ là đỉnh cuối cùng bị đưa lên gốc cây phụ trợ, nên
$5$ là tổ tiên chung gần nhất của $6$ và $9$.

\subsection{\texttt{findroot}}

\texttt{findroot} tìm gốc của cây gốc, không phải gốc của Link-Cut Tree.  Chính
xác hơn, sau khi dựng đường chính, nó tìm đỉnh có độ sâu nhỏ nhất trên đường
đó.

Sau \texttt{access(v)}, gốc cây gốc chắc chắn là đỉnh nhỏ nhất trong
Auxiliary Tree chứa $v$.  Ta đưa $v$ lên gốc Splay Tree, rồi đi liên tục sang
con trái cho đến khi không đi được nữa.  Vì Splay Tree lưu đường theo thứ tự
độ sâu, đỉnh ngoài cùng bên trái chính là gốc cây gốc.

\begingroup
\small
\begin{verbatim}
int findroot(int root)
{
    access(root);
    splay(root);
    int temp = root;
    while (t[temp].ch[0]) temp = t[temp].ch[0];
    return temp;
}
\end{verbatim}
\endgroup

Từ đây trở đi, ta xem Splay Tree như một công cụ; điều cần quan sát là thao tác
vĩ mô trên các đường ưu tiên, không phải từng biến đổi nhỏ bên trong Splay.

\subsection{\texttt{makeroot} hay \texttt{changeroot}}

Nhiều thao tác cần đổi gốc của cây gốc.  Trước hết dùng \texttt{access} để đặt
gốc cũ và gốc mới vào cùng một đường chính.  Khi đổi gốc từ $1$ sang $6$,
các đường phụ khác không đổi; chỉ có đường chính bị đảo ngược.  Vì thế
\texttt{makeroot(x)} chỉ cần dựng đường từ gốc tới $x$, đưa $x$ lên gốc cây
phụ trợ, rồi đánh dấu lười đảo đoạn.

\begingroup
\small
\begin{verbatim}
void changeroot(int root)
{
    access(root);
    splay(root);
    t[root].lazy ^= 1;
}
\end{verbatim}
\endgroup

\subsection{\texttt{link}}

Theo định nghĩa, để nối $v$ vào $w$, có thể truy cập $v$, đổi path parent của
cây phụ trợ chứa $v$ thành $w$, rồi truy cập lại $v$.  Trong cài đặt, ta làm
ngắn hơn: biến \texttt{root1} thành gốc cây gốc, rồi đặt cha của nó là
\texttt{root2}.  Sau khi nối, gốc của cây hợp nhất là gốc của cây chứa
\texttt{root2}; nếu cần một gốc khác, ta gọi \texttt{makeroot}.

\begingroup
\small
\begin{verbatim}
void link(int root1, int root2)
{
    makeroot(root1);
    t[root1].fa = root2;
}
\end{verbatim}
\endgroup

\subsection{\texttt{cut}}

Nếu \texttt{root1} và \texttt{root2} đang kề nhau trong cây, ta biến
\texttt{root1} thành gốc, rồi \texttt{access(root2)}.  Khi đó đường chính chỉ
gồm hai đỉnh \texttt{root1} và \texttt{root2}; sau \texttt{splay(root2)},
\texttt{root1} là con trái của \texttt{root2}.  Cắt cạnh bằng cách xóa con
trái ấy và xóa cha của \texttt{root1}.

\begingroup
\small
\begin{verbatim}
void cut(int root1, int root2)
{
    makeroot(root1);
    access(root2);
    splay(root2);
    t[root2].ch[0] = t[root1].fa = 0;
}
\end{verbatim}
\endgroup

\section{Cách viết Splay Tree trong Link-Cut Tree}

Splay Tree dùng trong Link-Cut Tree hơi khác Splay Tree thông thường.

\subsection{\texttt{isroot}}

Trong Splay Tree độc lập, một đỉnh là gốc nếu cha của nó bằng $0$.  Trong
Link-Cut Tree, gốc của một cây phụ trợ vẫn có thể trỏ cha tới path parent, tức
một đỉnh thuộc cây phụ trợ khác.  Vì cây phụ trợ kia không trỏ ngược lại bằng
con trái hoặc con phải, ta kiểm tra một đỉnh có phải gốc Splay Tree hay không
bằng cách xem cha của nó có nhận nó làm con hay không.

\begingroup
\small
\begin{verbatim}
bool isroot(int root)
{
    return t[t[root].fa].ch[0] != root
        && t[t[root].fa].ch[1] != root;
}
\end{verbatim}
\endgroup

\subsection{Xoay}

Khi xoay, cần chú ý ông của đỉnh có nằm trong cùng Splay Tree hay không.  Nếu
không, không được đặt con của ông thành đỉnh đang xoay; chỉ cần đặt cha của
đỉnh đang xoay thành ông.  So với Splay thường, khác biệt chính là kiểm tra
\texttt{isroot(fa)}.

\begingroup
\small
\begin{verbatim}
void rot(int root)
{
    int fa = t[root].fa;
    int gfa = t[fa].fa;
    int t1 = (root != t[fa].ch[0]);
    int t2 = (fa != t[gfa].ch[0]);
    int ch = t[root].ch[1 ^ t1];

    if (!isroot(fa)) t[gfa].ch[0 ^ t2] = root;

    t[ch].fa = fa;
    t[root].fa = gfa;
    t[root].ch[1 ^ t1] = fa;
    t[fa].fa = root;
    t[fa].ch[0 ^ t1] = ch;
    /* If extra data is maintained, update(fa) here. */
}
\end{verbatim}
\endgroup

\subsection{\texttt{splay}}

Khi \texttt{splay}, thứ nhất không được xoay sang cây phụ trợ khác.  Thứ hai,
các đánh dấu lười phải được đẩy xuống theo đúng thứ tự từ trên xuống dưới.  Ta
đi từ đỉnh hiện tại lên gốc Splay Tree, lưu đường đi vào một ngăn xếp, rồi
\texttt{pushdown} ngược lại.

\begingroup
\small
\begin{verbatim}
void pushdown(int root)
{
    if (t[root].lazy)
    {
        t[root].lazy = 0;
        swap(t[root].ch[0], t[root].ch[1]);
        t[t[root].ch[0]].lazy ^= 1;
        t[t[root].ch[1]].lazy ^= 1;
    }
}

void splay(int root)
{
    stk[++top] = root;
    for (int i = root; !isroot(i); i = t[i].fa)
        stk[++top] = t[i].fa;

    while (top) pushdown(stk[top--]);

    while (!isroot(root))
    {
        int fa = t[root].fa;
        int gfa = t[fa].fa;
        if (!isroot(fa))
            root == t[fa].ch[0] ^ fa == t[gfa].ch[0]
                ? rot(root) : rot(fa);
        rot(root);
    }
}
\end{verbatim}
\endgroup

\section{Bài tập thực hành}

\subsection{Duy trì động cấu trúc cây: Cave Survey}

Bài \textit{Cave Survey} (codevs1839) cho $n$ hang động và $m$ lệnh.  Ban đầu
không có đường hầm.  Mỗi lệnh là:
\begin{itemize}
  \item \texttt{Connect u v}: xuất hiện một đường hầm giữa $u$ và $v$;
  \item \texttt{Destroy u v}: đường hầm giữa $u$ và $v$ bị phá;
  \item \texttt{Query u v}: hỏi $u$ và $v$ có liên thông hay không.
\end{itemize}

Đề bảo đảm tại mọi thời điểm giữa hai hang bất kỳ có nhiều nhất một đường đi,
tức đồ thị luôn là một rừng.  Vì có xóa cạnh nên DSU thường không dùng trực
tiếp được.  Với Link-Cut Tree, kiểm tra liên thông bằng cách so sánh gốc của
hai đỉnh: nếu \texttt{findroot(u) == findroot(v)} thì liên thông.

Dữ liệu: $n\le 10000$, $m\le 200000$.  Mọi lệnh \texttt{Destroy} đều phá một
cạnh đang tồn tại.  Do dữ liệu lớn, cài đặt C/C++ nên dùng \texttt{scanf} và
\texttt{printf}.

\begingroup
\small
\begin{verbatim}
const int MAXN = 100005;
struct tree { int fa, ch[2], lazy; } t[MAXN];
int stk[MAXN], top;

void pushdown(int root) {
    if (t[root].lazy) {
        t[root].lazy = 0;
        swap(t[root].ch[0], t[root].ch[1]);
        t[t[root].ch[0]].lazy ^= 1;
        t[t[root].ch[1]].lazy ^= 1;
    }
}

bool isroot(int root) {
    return t[t[root].fa].ch[0] != root
        && t[t[root].fa].ch[1] != root;
}

void rot(int root) {
    int fa = t[root].fa, gfa = t[fa].fa;
    int t1 = (root != t[fa].ch[0]);
    int t2 = (fa != t[gfa].ch[0]);
    int ch = t[root].ch[1 ^ t1];
    if (!isroot(fa)) t[gfa].ch[0 ^ t2] = root;
    t[ch].fa = fa;
    t[root].fa = gfa;
    t[root].ch[1 ^ t1] = fa;
    t[fa].fa = root;
    t[fa].ch[0 ^ t1] = ch;
}

void splay(int root) {
    stk[++top] = root;
    for (int i = root; !isroot(i); i = t[i].fa)
        stk[++top] = t[i].fa;
    while (top) pushdown(stk[top--]);
    while (!isroot(root)) {
        int fa = t[root].fa, gfa = t[fa].fa;
        if (!isroot(fa))
            root == t[fa].ch[0] ^ fa == t[gfa].ch[0]
                ? rot(root) : rot(fa);
        rot(root);
    }
}

void access(int root) {
    int temp = 0;
    while (root) {
        splay(root);
        t[root].ch[1] = temp;
        temp = root;
        root = t[root].fa;
    }
}

void makeroot(int root) {
    access(root);
    splay(root);
    t[root].lazy ^= 1;
}

void link(int root1, int root2) {
    makeroot(root1);
    t[root1].fa = root2;
}

void cut(int root1, int root2) {
    makeroot(root1);
    access(root2);
    splay(root2);
    t[root2].ch[0] = t[root1].fa = 0;
}

int findroot(int root) {
    access(root);
    splay(root);
    int temp = root;
    while (t[temp].ch[0]) temp = t[temp].ch[0];
    return temp;
}

/* main: read commands Q/C/D.
   Q: print Yes if findroot(a) == findroot(b), otherwise No.
   C: link(a,b).
   D: cut(a,b). */
\end{verbatim}
\endgroup

\subsection{Duy trì bài toán đường đi đơn giản: Bouncing Sheep}

Bài \textit{Bouncing Sheep} (codevs2333) có $n$ thiết bị bật, đánh số
$0$ đến $n-1$.  Tại vị trí $i$, thiết bị đẩy con cừu tới $i+k_i$; nếu vượt ra
ngoài thì con cừu bay ra khỏi dãy.  Hai thao tác:
\begin{itemize}
  \item hỏi từ thiết bị $j$ cần bao nhiêu lần bật để bay ra;
  \item đổi hệ số bật của thiết bị $j$ thành $k$.
\end{itemize}

Thêm một đỉnh giả $n+1$ biểu diễn trạng thái đã bay ra.  Mỗi thiết bị có đúng
một cạnh đi tới thiết bị kế tiếp hoặc đỉnh giả; vì thế ta có một rừng động.
Đáp án truy vấn là số đỉnh trên đường từ $j$ tới $n+1$ trừ $1$.  Trong Splay
Tree cần duy trì kích thước đoạn đường.

\begingroup
\small
\begin{verbatim}
struct tree { int fa, ch[2], size, lazy; } t[MAXN];
int nxt[MAXN];

int nexts(int x, int y) {
    if (x + y > n) return n + 1;
    return x + y;
}

void updata(int root) {
    t[root].size = t[t[root].ch[0]].size
                 + t[t[root].ch[1]].size + 1;
}

void work(int root1, int root2) {
    makeroot(root1);
    access(root2);
    splay(root2);
}

/* Init: for every i, nxt[i] = nexts(i, k_i), link(i, nxt[i]).
   Type-1 query at a:
       ++a; work(a, n + 1); print t[n + 1].size - 1.
   Type-2 update at a with b:
       ++a; temp = nexts(a, b);
       if temp != nxt[a]: cut(a, nxt[a]); nxt[a] = temp; link(a, nxt[a]).
*/
\end{verbatim}
\endgroup

Tài liệu gốc nhấn mạnh một tối ưu nhỏ: giảm số lần gọi \texttt{update} có thể
tăng tốc đáng kể ở những bài phải duy trì nhiều thông tin trên Splay Tree.  Với
bài này hiệu quả không quá rõ, nhưng ở các bài nặng hơn thì đáng kể.

Các lỗi thường gặp:
\begin{enumerate}
  \item Vòng lặp vô hạn: phần lớn do viết sai \texttt{splay}.  Khi gỡ lỗi, in
  \texttt{root}, cha và ông của nó trong mỗi bước.
  \item Kết quả sai: sau \texttt{work}, dùng hàm \texttt{debug} để kiểm tra
  \texttt{root1} có thật là gốc không, và đường từ \texttt{root2} tới
  \texttt{root1} có thật là đường chính không.  Nếu cấu trúc sai, kiểm tra
  \texttt{access}.  Nếu cấu trúc đúng, kiểm tra đánh dấu lười trên đường và
  hàm \texttt{update}.
\end{enumerate}

\subsection{Trọng số cạnh hay trọng số đỉnh}

Bài \textit{Small Computer Room Tree} (codevs2370) cho một cây $N$ đỉnh đánh
số từ $0$ đến $N-1$, mỗi cạnh có chi phí $c$.  Với mỗi truy vấn hai đỉnh
$u,v$, cần in độ dài đường đi ngắn nhất giữa chúng.  Đây là bài LCA khá trần,
nhưng dùng để minh họa cách Link-Cut Tree xử lý trọng số cạnh.

Link-Cut Tree thuận tiện duy trì thông tin trên đỉnh, không trực tiếp trên
cạnh.  Cách chuẩn là biến mỗi cạnh thành một đỉnh mới mang trọng số của cạnh
đó, rồi nối $u$ với đỉnh cạnh và đỉnh cạnh với $v$.  Khi hỏi đường đi, dựng
đường từ $u$ tới $v$ và lấy tổng trên Splay Tree.

\begingroup
\small
\begin{verbatim}
const int MAXN = 225005;
struct tree { int fa, ch[2], num, sum, lazy; } t[MAXN];

void updata(int root) {
    t[root].sum = t[t[root].ch[0]].sum
                + t[t[root].ch[1]].sum
                + t[root].num;
}

void work(int root1, int root2) {
    changeroot(root1);
    access(root2);
    splay(root2);
}

int main() {
    scanf("%d", &n);
    for (int i = 1; i < n; ++i) {
        scanf("%d %d %d", &a1, &b1, &t[n+i].num);
        ++a1; ++b1;
        t[n+i].sum = t[n+i].num;
        link(a1, n+i);
        link(n+i, b1);
    }
    scanf("%d", &m);
    for (int i = 1; i <= m; ++i) {
        scanf("%d %d", &a1, &b1);
        ++a1; ++b1;
        work(a1, b1);
        printf("%d\n", t[b1].sum);
    }
}
\end{verbatim}
\endgroup

\subsection{Cây tĩnh với thông tin đoạn phức tạp hơn}

Bài \textit{Tree Statistics} (codevs2460, tuyển chọn đội tuyển tỉnh Chiết
Giang 2008) cho cây $n$ đỉnh, mỗi đỉnh có trọng số $w$.  Các thao tác:
\begin{itemize}
  \item \texttt{CHANGE u t}: đổi trọng số đỉnh $u$ thành $t$;
  \item \texttt{QMAX u v}: hỏi trọng số lớn nhất trên đường $u$--$v$;
  \item \texttt{QSUM u v}: hỏi tổng trọng số trên đường $u$--$v$.
\end{itemize}

Dữ liệu: $n\le 30000$, $q\le 200000$, trọng số nằm trong đoạn
$[-30000,30000]$.

Điểm cần giải thích là sửa một đỉnh: \texttt{splay(u)} để đưa đỉnh đó lên gốc
cây phụ trợ, sửa trực tiếp \texttt{num}, rồi \texttt{update}.  Với truy vấn
đường đi, dùng \texttt{makeroot(u); access(v); splay(v)} rồi đọc thông tin ở
\texttt{v}.

\begingroup
\small
\begin{verbatim}
struct tree { int ch[2], num, fa, lazy, maxx, sums; } t[MAXN];

void updata(int root) {
    t[root].maxx = max(t[root].num,
        max(t[t[root].ch[0]].maxx, t[t[root].ch[1]].maxx));
    t[root].sums = t[t[root].ch[0]].sums
                 + t[t[root].ch[1]].sums
                 + t[root].num;
}

void work(int root1, int root2) {
    makeroot(root1);
    access(root2);
    splay(root2);
}

/* After linking the initial edges:
   t[0].sums = 0; t[0].maxx = -INF.
   CHANGE a b: splay(a); t[a].num = b; updata(a).
   QMAX a b: work(a,b); in t[b].maxx.
   QSUM a b: work(a,b); in t[b].sums.
*/
\end{verbatim}
\endgroup

\subsection{Dùng nhiều đánh dấu lười và \texttt{update} phức tạp}

Bài \textit{Coloring} (codevs1566, tuyển chọn đội tuyển tỉnh Sơn Đông) cho cây
vô hướng $n$ đỉnh và $m$ thao tác:
\begin{itemize}
  \item \texttt{C a b c}: tô toàn bộ các đỉnh trên đường $a$--$b$ thành màu
  $c$;
  \item \texttt{Q a b}: hỏi số đoạn màu trên đường $a$--$b$, trong đó các đỉnh
  liên tiếp cùng màu được xem là một đoạn.
\end{itemize}

Bài này có sửa đoạn, nên cần đánh dấu lười.  Đồng thời tồn tại hai loại đánh
dấu: đánh dấu đảo chiều và đánh dấu gán màu.  Hai loại này không cản trở nhau
trong bài này.  Nếu là đánh dấu gán và đánh dấu cộng thì chúng có thể ảnh
hưởng lẫn nhau, khi đó phải xử lý thứ tự cẩn thận hơn.

Khi tính \texttt{update}, nếu con còn đánh dấu lười chưa đẩy xuống thì kết quả
có thể sai, tương tự truy vấn cực trị trong cây đoạn.  Vì vậy trước khi tổng
hợp, mã nguồn đẩy đánh dấu xuống hai con.  Nếu mọi giá trị nhỏ nhất trả về
thành $0$, cần kiểm tra đỉnh ảo $0$ đã được gán giá trị biên phù hợp chưa.

\begingroup
\small
\begin{verbatim}
struct tree {
    int ch[2], fa, lx, rx, num, col, lazy1, lazy2;
} t[MAXN];

void pushdown(int root) {
    if (t[root].lazy1) {
        swap(t[root].ch[0], t[root].ch[1]);
        swap(t[root].lx, t[root].rx);
        t[root].lazy1 = 0;
        t[t[root].ch[0]].lazy1 ^= 1;
        t[t[root].ch[1]].lazy1 ^= 1;
    }
    if (t[root].lazy2) {
        int c = t[root].lazy2;
        t[t[root].ch[0]].lazy2 = c;
        t[t[root].ch[1]].lazy2 = c;
        t[root].lx = t[root].rx = t[root].col = c;
        t[root].num = 1;
        t[root].lazy2 = 0;
    }
}

void updata(int root) {
    if (t[root].ch[0]) pushdown(t[root].ch[0]);
    if (t[root].ch[1]) pushdown(t[root].ch[1]);
    t[root].num = t[t[root].ch[0]].num
                + t[t[root].ch[1]].num + 1;
    if (t[t[root].ch[0]].rx == t[root].col) --t[root].num;
    if (t[t[root].ch[1]].lx == t[root].col) --t[root].num;
    t[root].lx = t[root].ch[0] ? t[t[root].ch[0]].lx : t[root].col;
    t[root].rx = t[root].ch[1] ? t[t[root].ch[1]].rx : t[root].col;
}

void change(int root1, int root2, int cols) {
    makeroot(root1);
    access(root2);
    splay(root2);
    t[root2].lazy2 = cols;
}

void work(int root1, int root2) {
    makeroot(root1);
    access(root2);
    splay(root2);
}

/* Q a b: work(a,b); in t[b].num.
   C a b c: change(a,b,c). */
\end{verbatim}
\endgroup

Tác giả gốc nhận xét rằng bài này khó viết đúng ngay lần đầu; nó thích hợp để
luyện kỹ năng gỡ lỗi Link-Cut Tree.

\subsection{Fruit Street III}

Bài \textit{Fruit Street III} (codevs3306) cũng minh họa tác động của đánh dấu
đảo chiều lên thông tin cần duy trì.  Có $n$ cửa hàng trên một cây, mỗi cửa
hàng có giá táo.  Hai thao tác:
\begin{itemize}
  \item \texttt{0 x y}: đổi giá cửa hàng $x$ thành $y$;
  \item \texttt{1 x y}: trên đường từ $x$ tới $y$, chọn một cửa hàng để mua và
  một cửa hàng về sau trên đường để bán, hỏi lợi nhuận lớn nhất.
\end{itemize}

Cần duy trì trên mỗi Splay Tree:
\begin{itemize}
  \item \texttt{maxx}: giá lớn nhất;
  \item \texttt{mins}: giá nhỏ nhất;
  \item \texttt{ans1}: lợi nhuận tốt nhất theo hướng trái sang phải;
  \item \texttt{ans2}: lợi nhuận tốt nhất theo hướng phải sang trái.
\end{itemize}

Khi đảo đoạn, hai con bị đổi chỗ và hai đáp án có hướng \texttt{ans1},
\texttt{ans2} cũng phải đổi chỗ.

\begingroup
\small
\begin{verbatim}
struct tree {
    int ch[2], fa, maxx, mins, num, ans1, ans2, lazy;
} t[MAXN];

void pushdown(int root) {
    if (t[root].lazy) {
        t[root].lazy = 0;
        swap(t[root].ch[0], t[root].ch[1]);
        swap(t[root].ans1, t[root].ans2);
        t[t[root].ch[0]].lazy ^= 1;
        t[t[root].ch[1]].lazy ^= 1;
    }
}

void update(int root) {
    if (t[root].ch[0]) pushdown(t[root].ch[0]);
    if (t[root].ch[1]) pushdown(t[root].ch[1]);
    int l = t[root].ch[0], r = t[root].ch[1];
    t[root].maxx = max(t[root].num, max(t[l].maxx, t[r].maxx));
    t[root].mins = min(t[root].num, min(t[l].mins, t[r].mins));
    t[root].ans1 = max(max(t[r].maxx, t[root].num)
                     - min(t[l].mins, t[root].num),
                       max(t[l].ans1, t[r].ans1));
    t[root].ans2 = max(max(t[l].maxx, t[root].num)
                     - min(t[r].mins, t[root].num),
                       max(t[l].ans2, t[r].ans2));
}

void change(int root, int d) {
    splay(root);
    t[root].num = d;
    update(root);
}

/* Query 1 x y: work(x,y); print t[y].ans1. */
\end{verbatim}
\endgroup

\subsection{Cây khung nhỏ nhất động: Comfortable Route}

Bài \textit{Comfortable Route} (codevs1001) cho $N$ điểm du lịch và $M$ con
đường hai chiều, mỗi đường có tốc độ bắt buộc $v_i$.  Với hai điểm $s,t$, cần
tìm một đường đi sao cho tỉ số giữa tốc độ lớn nhất và nhỏ nhất trên đường là
nhỏ nhất; nếu không có đường đi thì in \texttt{IMPOSSIBLE}, nếu cần thì in
phân số tối giản.

\subsubsection*{Cách DSU đơn giản}

Sắp xếp cạnh theo trọng số tăng dần.  Lần lượt chọn mỗi cạnh làm
\texttt{ans\_max}; sau đó duyệt ngược các cạnh không lớn hơn nó và dùng DSU để
nối các thành phần liên thông.  Khi $s$ và $t$ lần đầu liên thông, cạnh hiện
tại chính là \texttt{ans\_min} lớn nhất có thể ứng với \texttt{ans\_max} này.
Cập nhật tỉ số tốt nhất.

Nguyên lý là tham lam theo \texttt{ans\_min}: khi đã cố định tốc độ lớn nhất,
tốc độ nhỏ nhất càng lớn càng tốt.  Độ phức tạp theo mã gốc là
$O(n^2\log n)$, bộ nhớ $O(n)$.

\begingroup
\small
\begin{verbatim}
sort(edge + 1, edge + 1 + m, by_weight);
for (int i = 1; i <= m; ++i) {
    reset_dsu();
    for (int j = i; j >= 1; --j) {
        unite(edge[j].u, edge[j].v);
        if (find(s) == find(t)) {
            relax(edge[i].w, edge[j].w);
            break;
        }
    }
}
\end{verbatim}
\endgroup

\subsubsection*{Cách SPFA}

Một hướng khác là cố định giới hạn lớn nhất \texttt{max}.  Với các cạnh có
trọng số không vượt quá \texttt{max}, dùng SPFA để tối đa hóa giá trị nhỏ nhất
trên đường từ $s$ tới mọi đỉnh.  Nếu duyệt \texttt{max} tăng dần, mỗi lần chỉ
cần đưa hai đầu của cạnh mới vào hàng đợi và không cần xóa mảng \texttt{dis}.
Khi \texttt{dis[t]} tăng, \texttt{ans\_max} là trọng số đang duyệt và
\texttt{ans\_min} là \texttt{dis[t]}.

\begingroup
\small
\begin{verbatim}
int spfa(int maxvi)
{
    while (head != tail)
    {
        int x = pop_queue();
        for (edge i from x)
            if (i.vi <= maxvi
             && min(dmin[x], i.vi) > dmin[i.to])
            {
                dmin[i.to] = min(dmin[x], i.vi);
                push(i.to);
            }
    }
    if (dmin[t] changed) relax(maxvi, dmin[t]);
}
\end{verbatim}
\endgroup

\subsubsection*{Cách Link-Cut Tree}

Cách dùng Link-Cut Tree tương tự bài \textit{Magic Forest}.  Duyệt các cạnh
theo trọng số tăng dần, coi cạnh đang xét là ứng viên cho trọng số lớn nhất.
Ta duy trì một cây khung cực đại theo các cạnh đã xét: trên đường giữa hai đỉnh
bất kỳ, cạnh nhỏ nhất trên cây khung này là lớn nhất có thể.

Khi thêm cạnh $(u,v,w)$:
\begin{itemize}
  \item nếu $u$ và $v$ chưa liên thông, thêm cạnh vào cây;
  \item nếu đã liên thông, cạnh mới tạo một chu trình.  Tìm cạnh có trọng số
  nhỏ nhất trên đường $u$--$v$; nếu nó nhỏ hơn $w$, cắt cạnh nhỏ nhất ấy và
  thêm cạnh mới, ngược lại bỏ qua cạnh mới.
\end{itemize}

Sau mỗi bước, nếu $s$ và $t$ liên thông, dựng đường $s$--$t$, lấy trọng số nhỏ
nhất trên đường làm \texttt{ans\_min}, còn cạnh đang duyệt làm
\texttt{ans\_max}, rồi cập nhật đáp án.  Độ phức tạp khấu hao là
$O(m\log n)$, bộ nhớ $O(n+m)$.

Như ở phần trọng số cạnh, mỗi cạnh được biểu diễn bằng một đỉnh mới mang trọng
số cạnh; các đỉnh gốc nhận trọng số vô cùng lớn để không ảnh hưởng tới phép
lấy min.

\begingroup
\small
\begin{verbatim}
struct tree { int fa, ch[2], num, mins, lazy; } t[MAXN];

void update(int root) {
    t[root].mins = root;
    if (t[t[t[root].ch[0]].mins].num < t[t[root].mins].num)
        t[root].mins = t[t[root].ch[0]].mins;
    if (t[t[t[root].ch[1]].mins].num < t[t[root].mins].num)
        t[root].mins = t[t[root].ch[1]].mins;
}

for (int i = 1; i <= m; ++i) {
    int u = e[i].p[0], v = e[i].p[1], w = e[i].vi;
    if (findroot(u) != findroot(v)) {
        t[i+n].num = w;
        t[i+n].mins = i+n;
        t[u].num = t[v].num = INF;
        link(u, i+n);
        link(i+n, v);
    } else {
        works(u, v);
        int old = t[v].mins;
        if (t[old].num < w) {
            cut(e[old-n].p[0], old);
            cut(old, e[old-n].p[1]);
            t[i+n].num = w;
            t[i+n].mins = i+n;
            link(u, i+n);
            link(i+n, v);
        }
    }

    if (findroot(s) == findroot(t)) {
        works(s, t);
        relax(w, t[t[t1].mins].num); /* idea: min edge on path s--t */
    }
}
\end{verbatim}
\endgroup

Dòng cuối trong bản trích trên chỉ diễn đạt ý tưởng của mã nguồn gốc: sau
\texttt{works(s,t)}, thông tin min nằm ở gốc Splay Tree của đường $s$--$t$.
Trong mã hoàn chỉnh, biến đích là \texttt{t1} và biểu thức tương ứng là trọng
số của đỉnh cạnh nhỏ nhất trên đường ấy.

Tác giả gốc kết luận bằng gợi ý: hãy dùng cả SPFA và Link-Cut Tree để giải
\textit{NOI Magic Forest}; bài đó gần như cùng khuôn mẫu với lời giải Link-Cut
Tree ở trên.

\section{Dữ liệu mẫu trong các bài thực hành}

Phần mã nguồn trong tài liệu gốc rất dài; theo ngoại lệ của dự án, bản dịch
không chép lại toàn bộ chương trình mẫu. Tuy nhiên các bộ dữ liệu mẫu của bài
tập được giữ lại để người đọc đối chiếu phát biểu bài.

\subsection*{Cave Survey}
\begin{verbatim}
Input
200 5
Query 123 127
Connect 123 127
Query 123 127
Destroy 127 123
Query 123 127

Output
No
Yes
No
\end{verbatim}

\subsection*{Bouncing Sheep}
\begin{verbatim}
Input
4
1 2 1 1
3
1 1
2 1 1
1 1

Output
2
3
\end{verbatim}

\subsection*{Small Computer Room Tree}
\begin{verbatim}
Input
3
1 0 1
2 0 1
3
1 0
2 0
1 2

Output
1
1
2
\end{verbatim}

\subsection*{Tree Statistics}
\begin{verbatim}
Input
4
1 2
2 3
4 1
4 2 1 3
12
QMAX 3 4
QMAX 3 3
QMAX 3 2
QMAX 2 3
QSUM 3 4
QSUM 2 1
CHANGE 1 5
QMAX 3 4
CHANGE 3 6
QMAX 3 4
QMAX 2 4
QSUM 3 4

Output
4
1
2
2
10
6
5
6
5
16
\end{verbatim}

\subsection*{Coloring}
\begin{verbatim}
Input
6 5
2 2 1 2 1 1
1 2
1 3
2 4
2 5
2 6
Q 3 5
C 2 1 1
Q 3 5
C 5 1 2
Q 3 5

Output
3
1
2
\end{verbatim}

\subsection*{Fruit Street III}
\begin{verbatim}
Input
4
16 5 1 15
1 2
1 3
2 4
3
1 3 4
0 3 17
1 4 3

Output
15
12
\end{verbatim}

\subsection*{Comfortable Route}
\begin{verbatim}
Sample 1 input
4 2
1 2 1
3 4 2
1 4

Sample 1 output
IMPOSSIBLE

Sample 2 input
3 3
1 2 10
1 2 5
2 3 8
1 3

Sample 2 output
5/4

Sample 3 input
3 2
1 2 2
2 3 4
1 3

Sample 3 output
2
\end{verbatim}

\section{Ghi chú về độ phức tạp}

Các thao tác cốt lõi \texttt{access}, \texttt{makeroot}, \texttt{link},
\texttt{cut}, \texttt{findroot} đều được xây trên Splay Tree.  Mỗi lần
\texttt{access} có thể thay đổi nhiều cạnh ưu tiên, nhưng phân tích khấu hao
của Sleator--Tarjan cho thấy tổng chi phí được chặn bởi $O(\log n)$ cho mỗi
thao tác.  Vì vậy, trong các bài cây động, mỗi cập nhật hoặc truy vấn đường đi
thường đạt $O(\log n)$ khấu hao, cộng thêm chi phí duy trì thông tin đoạn trong
\texttt{update} và \texttt{pushdown}.

\end{document}
